\documentclass[a4paper,landscape,8pt]{extarticle}
\usepackage[landscape,margin=0.7cm,top=0.5cm,bottom=0.7cm]{geometry}
\usepackage[utf8]{inputenc}
\usepackage[T1]{fontenc}
\usepackage[ngerman]{babel}
\usepackage{helvet}
\renewcommand{\familydefault}{\sfdefault}
\usepackage{xcolor}
\usepackage{tikz}
\usetikzlibrary{positioning}
\usepackage{enumitem}
\usepackage{multicol}
\usepackage{array}
\usepackage{fancyhdr}
\usepackage{amssymb}

% Farben
\definecolor{restartblue}{HTML}{1565C0}
\definecolor{wlangreen}{HTML}{43A047}
\definecolor{storageorange}{HTML}{FF9800}
\definecolor{lightgray}{HTML}{F5F5F5}
\definecolor{bordergray}{HTML}{BBBBBB}

\pagestyle{fancy}
\fancyhf{}
\renewcommand{\headrulewidth}{0pt}
\fancyfoot[C]{\footnotesize Seite \thepage{} von 3}

\setlength{\parindent}{0pt}
\setlength{\parskip}{2pt}

% Kompakte Listen
\setlist{nosep, leftmargin=1em, itemsep=0pt, topsep=1pt}

\begin{document}

% ============================================================================
% SEITE 1: INFORMATIONSBLATT
% ============================================================================

\noindent{\Large\bfseries Probleme selbst lösen: Die 5 Erste-Hilfe-Schritte} \hfill
{\footnotesize\textbf{Lernziele:} Neustart | App schließen | WLAN \& Speicher prüfen}

\vspace{1mm}
\hrule
\vspace{1mm}

% Drei-Spalten-Layout
\begin{minipage}[t]{0.32\linewidth}
\begin{center}
\colorbox{restartblue!15}{%
\begin{minipage}{0.95\linewidth}
\vspace{2mm}
\begin{center}
{\Large\bfseries\textcolor{restartblue}{1. Neustart \& App schließen}}\\[1pt]
{\footnotesize\itshape Die wichtigsten Tricks}
\end{center}
\vspace{2mm}
\end{minipage}%
}
\end{center}

\vspace{1mm}

\textbf{Goldene Regel:}\\
\textbf{Neustart löst 80\% aller Probleme!}

\textbf{iPhone/iPad neu starten:}
\begin{itemize}
\item \textbf{Mit Face ID:} Seitentaste + Lautstärke-Taste gleichzeitig drücken
\item ``Ausschalten'' nach rechts wischen
\item 10 Sekunden warten
\item Seitentaste drücken bis Apfel erscheint
\end{itemize}

\textbf{Mit Home-Button:}
\begin{itemize}
\item Seitentaste \textbf{lange} drücken
\item ``Ausschalten'' wischen
\item Wieder einschalten
\end{itemize}

\textbf{App schließen (wenn sie hängt):}
\begin{enumerate}
\item Von unten nach oben wischen
\item In der \textbf{Mitte stoppen}
\item Hängende App nach \textbf{oben wegwischen}
\item App neu öffnen
\end{enumerate}

\end{minipage}%
\hfill%
\begin{minipage}[t]{0.32\linewidth}
\begin{center}
\colorbox{wlangreen!15}{%
\begin{minipage}{0.95\linewidth}
\vspace{2mm}
\begin{center}
{\Large\bfseries\textcolor{wlangreen}{2. WLAN prüfen}}\\[1pt]
{\footnotesize\itshape Ist Internet da?}
\end{center}
\vspace{2mm}
\end{minipage}%
}
\end{center}

\vspace{1mm}

\textbf{Problem:} ``Das Internet geht nicht!''

\textbf{WLAN prüfen:}
\begin{enumerate}
\item Einstellungen öffnen
\item ``WLAN'' antippen
\item Ist der \textbf{Schalter grün}?
\item Steht ein \textbf{Häkchen} beim Netzwerk?
\end{enumerate}

\textbf{WLAN reparieren:}
\begin{itemize}
\item WLAN \textbf{aus- und wieder einschalten}
\item Netzwerk antippen $\rightarrow$ neu verbinden
\item Router neu starten (zu Hause)
\end{itemize}

\vspace{2mm}
\fcolorbox{wlangreen}{wlangreen!5}{%
\begin{minipage}{0.92\linewidth}
\footnotesize
\textbf{Tipp:} Kein WLAN-Symbol oben? Dann bist du im Mobilfunk -- das kostet evtl. Datenvolumen!
\end{minipage}%
}

\vspace{2mm}

\textbf{WLAN-Symbol:}\\
\begin{tikzpicture}[scale=0.5]
\draw[line width=1.5pt, wlangreen] (0,0) arc (180:0:0.3);
\draw[line width=1.5pt, wlangreen] (-0.15,0.2) arc (180:0:0.45);
\draw[line width=1.5pt, wlangreen] (-0.3,0.4) arc (180:0:0.6);
\fill[wlangreen] (0.3,-0.1) circle (0.08);
\end{tikzpicture}
= WLAN verbunden

\end{minipage}%
\hfill%
\begin{minipage}[t]{0.32\linewidth}
\begin{center}
\colorbox{storageorange!15}{%
\begin{minipage}{0.95\linewidth}
\vspace{2mm}
\begin{center}
{\Large\bfseries\textcolor{storageorange}{3. Speicher prüfen}}\\[1pt]
{\footnotesize\itshape Ist noch Platz?}
\end{center}
\vspace{2mm}
\end{minipage}%
}
\end{center}

\vspace{1mm}

\textbf{Problem:} ``Mein Handy ist so langsam!''

\textbf{Speicher prüfen:}
\begin{enumerate}
\item Einstellungen öffnen
\item ``Allgemein'' antippen
\item ``iPhone-Speicher''
\item Farbigen Balken anschauen
\end{enumerate}

\textbf{Was verbraucht Platz?}
\begin{itemize}
\item \textbf{Fotos/Videos} -- oft am meisten!
\item \textbf{Apps} -- besonders Spiele
\item \textbf{Nachrichten} -- mit vielen Bildern
\end{itemize}

\textbf{Speicher freigeben:}
\begin{itemize}
\item Alte Fotos löschen oder in iCloud
\item Ungenutzte Apps löschen
\item ``Zuletzt gelöscht'' leeren
\end{itemize}

\vspace{2mm}
\fcolorbox{storageorange}{storageorange!5}{%
\begin{minipage}{0.92\linewidth}
\footnotesize
\textbf{Faustregel:} Mind. 2-3 GB sollten frei sein, damit alles flüssig läuft.
\end{minipage}%
}

\end{minipage}

\vspace{1mm}

% Zusammenfassung
\begin{center}
\fcolorbox{bordergray}{lightgray}{%
\begin{minipage}{0.98\linewidth}
\centering
\vspace{1.5mm}
\textbf{Die 5 Erste-Hilfe-Schritte:}
1. \textcolor{restartblue}{\textbf{Ruhe bewahren}} -- Du machst nichts kaputt! ~$|$~
2. \textcolor{restartblue}{\textbf{Neustart}} ~$|$~
3. \textcolor{restartblue}{\textbf{App schließen}} ~$|$~
4. \textcolor{wlangreen}{\textbf{WLAN prüfen}} ~$|$~
5. \textcolor{storageorange}{\textbf{Speicher prüfen}}
\vspace{1.5mm}
\end{minipage}%
}
\end{center}

\newpage

% ============================================================================
% SEITE 2: ÜBUNGEN
% ============================================================================

\begin{center}
{\LARGE\bfseries Übungen: Probleme selbst lösen}\\[3pt]
{\normalsize Schau dir Seite 1 an und beantworte die Fragen}
\end{center}

\vspace{3mm}
\hrule
\vspace{4mm}

\begin{multicols}{2}

\textbf{\large Teil A: Die 5 Schritte}

Bringe die 5 Erste-Hilfe-Schritte in die richtige Reihenfolge (1--5):

\vspace{2mm}

\renewcommand{\arraystretch}{1.8}
\begin{tabular}{@{}cp{7cm}@{}}
$\square$ & Speicher prüfen \\
$\square$ & Ruhe bewahren \\
$\square$ & WLAN prüfen \\
$\square$ & Neustart \\
$\square$ & App schließen \\
\end{tabular}
\renewcommand{\arraystretch}{1}

\vspace{6mm}

\textbf{\large Teil B: Richtig oder Falsch?}

Kreuze an.

\vspace{2mm}

\renewcommand{\arraystretch}{1.6}
\begin{tabular}{@{}p{8cm}cc@{}}
& R & F \\
1. Neustart löst viele Probleme. & $\square$ & $\square$ \\
2. Eine hängende App muss man nach unten wischen. & $\square$ & $\square$ \\
3. WLAN ist aus, wenn der Schalter grün ist. & $\square$ & $\square$ \\
4. Fotos verbrauchen viel Speicherplatz. & $\square$ & $\square$ \\
5. ``Zuletzt gelöscht'' leeren spart Speicher. & $\square$ & $\square$ \\
\end{tabular}
\renewcommand{\arraystretch}{1}

\vspace{6mm}

\textbf{\large Teil C: Problem-Zuordnung}

Welcher Schritt hilft? Verbinde:

\vspace{2mm}

\renewcommand{\arraystretch}{1.8}
\begin{tabular}{@{}p{4cm}p{4.5cm}@{}}
``Handy ist langsam'' & $\circ$ ~~~~ $\circ$ Neustart \\
``App reagiert nicht'' & $\circ$ ~~~~ $\circ$ WLAN prüfen \\
``Seite lädt nicht'' & $\circ$ ~~~~ $\circ$ Speicher prüfen \\
``Alles hängt'' & $\circ$ ~~~~ $\circ$ App schließen \\
\end{tabular}
\renewcommand{\arraystretch}{1}

\columnbreak

\textbf{\large Teil D: Praktisch üben}

Führe diese Aufgaben auf deinem Gerät durch:

\vspace{3mm}

\textbf{Aufgabe 1: Speicher prüfen}\\
Gehe zu Einstellungen $\rightarrow$ Allgemein $\rightarrow$ iPhone-Speicher

\vspace{2mm}
Wie viel Speicher ist belegt? \dotfill

\vspace{2mm}
Was verbraucht am meisten? \dotfill

\vspace{2mm}
$\square$ Erledigt

\vspace{5mm}

\textbf{Aufgabe 2: WLAN prüfen}\\
Gehe zu Einstellungen $\rightarrow$ WLAN

\vspace{2mm}
Ist WLAN aktiviert (grün)? ~~ $\square$ Ja ~~ $\square$ Nein

\vspace{2mm}
Mit welchem Netzwerk bist du verbunden? \dotfill

\vspace{2mm}
$\square$ Erledigt

\vspace{5mm}

\textbf{Aufgabe 3: App-Umschalter öffnen}\\
Wische von unten nach oben und halte in der Mitte an.

\vspace{2mm}
Wie viele Apps sind offen? \dotfill

\vspace{2mm}
$\square$ Erledigt

\vspace{5mm}

\textbf{Bonus: Neustart durchführen}\\
Starte dein Gerät einmal komplett neu.

\vspace{2mm}
$\square$ Erledigt

\end{multicols}

\newpage

% ============================================================================
% SEITE 3: CHECKLISTE
% ============================================================================

\begin{center}
{\LARGE\bfseries Checkliste: Das kann ich jetzt}\\[3pt]
{\normalsize Geh die Liste durch und hake ab, was du sicher kannst}
\end{center}

\vspace{3mm}
\hrule
\vspace{6mm}

% Drei Boxen nebeneinander
\begin{minipage}[t]{0.31\linewidth}
\begin{center}
\fcolorbox{restartblue}{restartblue!10}{%
\begin{minipage}{0.92\linewidth}
\vspace{3mm}
\begin{center}
{\large\bfseries\textcolor{restartblue}{Neustart \& App}}
\end{center}
\vspace{2mm}
\begin{itemize}[label=$\square$, itemsep=5pt]
\item Ich kann mein Gerät \textbf{neu starten}
\item Ich kenne den \textbf{Unterschied} Face ID / Home-Button
\item Ich kann den \textbf{App-Umschalter} öffnen
\item Ich kann eine App \textbf{schließen} (wegwischen)
\end{itemize}
\vspace{3mm}
\end{minipage}%
}
\end{center}
\end{minipage}%
\hfill%
\begin{minipage}[t]{0.31\linewidth}
\begin{center}
\fcolorbox{wlangreen}{wlangreen!10}{%
\begin{minipage}{0.92\linewidth}
\vspace{3mm}
\begin{center}
{\large\bfseries\textcolor{wlangreen}{WLAN prüfen}}
\end{center}
\vspace{2mm}
\begin{itemize}[label=$\square$, itemsep=5pt]
\item Ich finde die \textbf{WLAN-Einstellung}
\item Ich erkenne, ob WLAN \textbf{an} ist (grün)
\item Ich sehe das \textbf{Häkchen} beim Netzwerk
\item Ich kann WLAN \textbf{aus- und einschalten}
\end{itemize}
\vspace{3mm}
\end{minipage}%
}
\end{center}
\end{minipage}%
\hfill%
\begin{minipage}[t]{0.31\linewidth}
\begin{center}
\fcolorbox{storageorange}{storageorange!10}{%
\begin{minipage}{0.92\linewidth}
\vspace{3mm}
\begin{center}
{\large\bfseries\textcolor{storageorange}{Speicher prüfen}}
\end{center}
\vspace{2mm}
\begin{itemize}[label=$\square$, itemsep=5pt]
\item Ich finde den \textbf{iPhone-Speicher}
\item Ich verstehe den \textbf{farbigen Balken}
\item Ich weiß, was \textbf{Speicher frisst}
\item Ich kenne Wege zum \textbf{Aufräumen}
\end{itemize}
\vspace{3mm}
\end{minipage}%
}
\end{center}
\end{minipage}

\vspace{8mm}

% Gesamtverständnis
\begin{center}
\fcolorbox{black}{lightgray}{%
\begin{minipage}{0.95\linewidth}
\vspace{4mm}
\begin{center}
{\large\bfseries Das Wichtigste}
\end{center}
\vspace{3mm}
\begin{multicols}{2}
\begin{itemize}[label=$\square$, itemsep=6pt]
\item Ich bleibe \textbf{ruhig} bei Problemen
\item Ich probiere \textbf{erst selbst}, bevor ich frage
\item Ich kenne die \textbf{5 Erste-Hilfe-Schritte}
\item Ich weiß: \textbf{Neustart hilft oft!}
\item Ich fühle mich \textbf{sicherer} im Umgang
\end{itemize}
\end{multicols}
\vspace{3mm}
\end{minipage}%
}
\end{center}

\vspace{8mm}

% Notfall-Notizen
\begin{center}
\fcolorbox{bordergray}{white}{%
\begin{minipage}{0.95\linewidth}
\vspace{4mm}
\begin{center}
{\large\bfseries Meine Notizen}\\[3pt]
{\footnotesize (Trage ein, was du dir merken möchtest)}
\end{center}
\vspace{4mm}
\renewcommand{\arraystretch}{2}
\begin{tabular}{@{}p{0.45\linewidth}p{0.45\linewidth}@{}}
\textbf{Mein WLAN-Name:} & \textbf{Mein Speicherstand:} \\
\dotfill & Belegt: \dotfill von \dotfill GB \\[2mm]
\textbf{Was bei mir am meisten Platz braucht:} & \textbf{Was ich aufräumen könnte:} \\
\dotfill & \dotfill \\[2mm]
\end{tabular}
\renewcommand{\arraystretch}{1}
\vspace{4mm}
\end{minipage}%
}
\end{center}

\vfill

\begin{center}
\footnotesize\textit{Stand: \today{} -- Erst selbst probieren, dann fragen. Neustart löst 80\% aller Probleme!}
\end{center}

\end{document}
